\documentclass[11pt,a4paper]{article}

\usepackage[margin=2.5cm]{geometry}
\usepackage{amsmath,amssymb}
\usepackage{booktabs}
\usepackage{array}
\usepackage{graphicx}
\usepackage{xcolor}
\usepackage[hidelinks]{hyperref}
\usepackage{microtype}
\usepackage{titlesec}
\usepackage{enumitem}
\usepackage{caption}
\usepackage{float}

\definecolor{accent}{HTML}{8a6d1f}
\titleformat{\section}{\Large\bfseries\color{accent}}{\thesection}{0.6em}{}
\titleformat{\subsection}{\large\bfseries}{\thesubsection}{0.6em}{}

\newcommand{\pct}{\,\%}
\newcommand{\code}[1]{\texttt{#1}}

\title{\bfseries meshprod: Monocular Video to Measured Mesh\\[0.3em]
\large A Technical Report on the Sparse SfM--Dense MVS--Fusion Pipeline}
\author{Adarsh Gupta\\Indian Institute of Information Technology, Gwalior (IIIT Gwalior)}
\date{\today}

\begin{document}

\maketitle

\begin{abstract}
\noindent
This report documents \code{meshprod}, a self-contained pipeline that turns a single
hand-held video of an object into a measured triangle mesh, with a machine-readable
assessment of how trustworthy the result is. The system chains nine classical
computer-vision stages---decode, feature extraction, matching, track building, incremental
structure-from-motion, dense multi-view stereo, TSDF fusion, marching-tetrahedra meshing and
simplification---on top of the \code{kornia-rs} library, and reports every number it emits.
We describe each stage with the mathematics and the research it derives from, present the
measured results on a set of reference captures, and catalogue the engineering iterations
that fixed concrete failures (an intrinsic-orientation bug, a floating-point cache
round-trip error, a fragile scale anchor, and a whole-scene fusion problem). The pipeline
contains no learned model; its correctness rests on measured reprojection and depth errors,
and on deterministic, reproducible behaviour.
\end{abstract}

\tableofcontents

\section{Introduction}
\label{sec:intro}

\code{meshprod} reconstructs a \emph{single object} filmed by a hand-held or orbiting camera,
producing a triangle mesh plus a machine-readable report that states how reliable the result
is. Two properties are non-negotiable:

\begin{itemize}[leftmargin=1.5em]
\item \textbf{Correctness over decoration.} Every number in the report is measured
(reprojection error, depth-vs-sparse disagreement, mesh topology); nothing is asserted
without a check behind it.
\item \textbf{Honest failure.} When the input cannot support a mesh, the pipeline stops with a
named error and actionable hints rather than emitting a confident blob.
\end{itemize}

The project is a thin application crate (\code{sfm2mesh}) built on the \code{kornia-rs}
Rust computer-vision stack, pinned to the \code{semproj/sfm2mesh} branch. The application owns
the orchestration, the dense/TSDF/mesh stages and a local web UI; \code{kornia-rs} provides
features, matching primitives, multi-view geometry, PnP, bundle adjustment and incremental SfM.

\section{System architecture}
\label{sec:arch}

\begin{center}\small
\begin{tabular}{lll}
\toprule
\textbf{Layer} & \textbf{Crate} & \textbf{Role}\\
\midrule
Geometry / SfM & \code{kornia-calib}, \code{kornia-3d} & tracks, PnP, BA, triangulation, pose-graph\\
Features / images & \code{kornia-imgproc}, \code{kornia-image} & SIFT/ORB, colour conversion\\
I/O & \code{kornia-io} & GStreamer decode, JPEG/PNG\\
Application & \code{sfm2mesh} & orchestration, dense, TSDF, mesh, UI\\
\bottomrule
\end{tabular}
\end{center}

The data flow is:

\begin{center}\small
\texttt{video}\ $\to$\ \texttt{decode}\ $\to$\ \texttt{features}\ $\to$\ \texttt{match}\ $\to$\ \texttt{tracks}\ $\to$\ \texttt{sparse SfM}\ $\to$\ [\texttt{scale}]\ $\to$\ \texttt{isolation}\\[0.25em]
$\to$\ \texttt{dense depth}\ $\to$\ \texttt{TSDF}\ $\to$\ \texttt{mesh}\ $\to$\ \texttt{report}
\end{center}

A local web server (\code{--serve}) runs the CLI as a child process, parses its progress and
metric lines, and serves a browser front-end plus the PLY artifacts.

\section{The pipeline}
\label{sec:pipeline}

\subsection{Decode}
\label{sec:decode}

The video is read through GStreamer into RGB frames, each converted to grayscale for feature
detection. Only every $N$-th frame is kept (\code{--frame-step}): consecutive 30\,fps frames
move too little to triangulate stably, so a step of $5$--$20$ sets a healthy baseline while
cutting decode and matching cost. Frames are cached because decode is the single most
expensive fixed cost. The \emph{intrinsics} must match the decoded frame, not the sensor: a
portrait clip decoded unrotated has its principal point swapped relative to a landscape one,
and $f_x \approx 0.7\times\max(w,h)$ is a good phone guess.

\subsection{Feature extraction}
\label{sec:features}

SIFT \cite{lowe2004} is the default: scale- and rotation-invariant, robust to moderate
viewpoint change, and sub-pixel localised. ORB \cite{rublee2011} is available as a faster,
weaker alternative. A capture with fewer than $\sim$80 keypoints per frame is almost always
blurry, defocused or textureless; the pipeline warns, because no downstream stage can recover
structure that was never detected.

\subsection{Matching}
\label{sec:matching}

Each frame is compared to the next $W$ frames (\code{--match-window}) using \emph{mutual
nearest neighbours} plus \emph{Lowe's ratio test} ($d_1 < r\,d_2$) \cite{lowe2004}. Optional
\emph{epipolar verification} fits a fundamental matrix per pair with RANSAC
\cite{fischler1981} and keeps geometrically consistent inliers; the seven-point estimator
needs at least seven points, and a pair with fewer than \code{--geo-min-inliers} inliers is
dropped \cite{hartley2004}. \emph{Wide-baseline} matching additionally links frame $i$ to
$i+K, i+2K, \ldots$ to bridge fast motion.

\subsection{Tracks}
\label{sec:tracks}

Matches are chained into multi-view tracks with a union--find over $(camera, keypoint)$ nodes
\cite{schonberger2016}. Two cases are handled explicitly: \emph{orientation siblings} (SIFT can
emit two keypoints at the same pixel, which are merged) and \emph{genuine ambiguity} (two
distinct image locations from one camera in one component, which drops the track). Track order
is made deterministic because floating-point summation in bundle adjustment depends on
observation order.

\subsection{Sparse structure-from-motion}
\label{sec:sfm}

\code{kornia\_calib::reconstruct} implements incremental SfM in the COLMAP tradition
\cite{schonberger2016}:

\begin{enumerate}[leftmargin=1.5em]
\item \textbf{Bootstrap} --- pick a seed pair by geometry (shared tracks, median triangulation
angle, homography-vs-fundamental ratio) and initialise a two-view essential solve
\cite{hartley2004}.
\item \textbf{Grow} --- register the unplaced camera with the most 2D--3D links via PnP-RANSAC;
EPnP \cite{lepetit2009} gives the pose, RANSAC rejects outliers, and SPRT \cite{wald1945}
streams residuals to reject hopeless hypotheses early. A view is accepted only with at least
\code{--min-registration-inliers} inliers.
\item \textbf{Triangulate} --- new tracks are triangulated with parallax, cheirality and
reprojection gates.
\item \textbf{Bundle adjustment} --- minimise reprojection error over poses and points with
Levenberg--Marquardt and a robust (Huber/Cauchy) loss via IRLS
\cite{triggs2000,lourakis2009}. The reduced camera system is factorised sparsely, which scales
linearly rather than cubically in the camera count.
\item \textbf{Maintenance} --- periodic global BA, reprojection filtering, re-admission of
previously rejected sightings (\emph{CompleteTracks}), and local windowed BA after each
registration keep long chains from drifting.
\end{enumerate}

The solve is \emph{up to scale} until something provides a metric reference.

\subsection{Calibration and metric scale}
\label{sec:scale}

Two distinct problems are conflated in the field and kept separate here:

\begin{itemize}[leftmargin=1.5em]
\item \textbf{Focal length} (intrinsics), estimated from a planar AprilTag \cite{olson2011}
via Zhang's orthonormality constraint \cite{zhang2000}: \code{estimate\_focal} recovers $f$
from the board-to-image homography, taking the median over sampled frames and reporting the
robust (MAD) spread. Single-view estimation cannot separate the principal point or distortion;
multi-view Zhang (\code{calibrate\_multishot}) is available for that.
\item \textbf{Metric scale} (extrinsics gauge). A tag of known edge length \emph{in the scene}
turns a shape into a measured shape. Because kornia's built-in \code{tag\_scale} triangulates
the tag from a single widest-baseline pair (fragile on an orbit), the application instead
triangulates each corner by \emph{N-view DLT} over every registered view that saw the tag and
takes the median edge. A scene extent that is implausible for the tag size is flagged.
\end{itemize}

\subsection{Foreground isolation}
\label{sec:isolation}

A ``scattered'' mesh is usually the whole room, not the object. The capture centres the object,
so the application scores every sparse point by the fraction of its observations in which it
projects inside a central disk of the image. Well-supported \emph{and} central points are kept,
the largest 26-connected component is taken, and an implausibly tiny cluster is rejected. The
mask trims the sparse cloud and gives the dense stage a tight box. Crucially, the TSDF
resolution is rescaled to the object box so the voxel size stays at the full-scene scale: a
tight box at the configured resolution over-resolves the depth noise and shatters the surface
(measured: 845 patches at $0.0004$\,units/voxel vs.\ 1 patch at $0.0032$).

\subsection{Dense depth}
\label{sec:dense}

Each registered view gets a depth map by \emph{plane sweeping} \cite{collins1996}: sample
inverse-depth hypotheses and score each by photometric consistency with neighbours. Costs
available: NCC (normalised cross-correlation), Census \cite{zabih1994} (intensity orderings),
and SGM \cite{hirschmuller2008} (Census plus a smoothness prior along 4 or 8 paths). SGM is the
default on low/medium-texture captures (measured: Census WTA produced 30 patches on the
keyboard clip, SGM produced 2). Winners are refined to sub-pixel by fitting a parabola, a
pixel is kept only if \code{--lr-required} neighbours agree, and a speckle filter removes
isolated pixels.

\subsection{TSDF fusion}
\label{sec:tsdf}

Depth maps are fused into a truncated signed distance field \cite{curless1996,newcombe2011}.
Each voxel accumulates confidence-weighted ``in front/behind'' votes; the volume is bounded by
the robust (1st--99th percentile) extent of the dense cloud, and the truncation band is
\code{--trunc-voxels} voxels wide. The truncation must exceed the depth noise or views disagree
and tear the surface.

\subsection{Meshing, smoothing, simplification}
\label{sec:mesh}

The zero level set is extracted with \emph{marching tetrahedra} (a variant of marching cubes
\cite{lorensen1987}). Small patches are dropped, the surface is smoothed with Taubin $\lambda/\mu$
passes \cite{taubin1995} and a bilateral pass that preserves creases, and finally simplified to
a face budget with quadric error metrics \cite{garland1997}. Colours and normals are assigned
per vertex by projecting into the source frames; normals come from the TSDF gradient.

\subsection{Report and named failures}
\label{sec:report}

Every run writes \code{<stem>\_report.json}: elapsed time, statistics (views, points,
reprojection RMSE, depth error percentiles, face counts), a mesh-quality pack (patches, Euler
characteristic, boundary-edge ratio, degenerate/needle fractions, closedness), provenance
(video SHA-256, exact argv, version), resources (peak RSS), and a heuristic \emph{self-assessment}
confidence. Failures write \code{<stem>\_error.json} with a stable machine code and hints.

\section{Results}
\label{sec:results}

All numbers were measured on one machine. Intrinsics were derived from the decoded frame
($f_x \approx 0.7 \times \max(w,h)$, principal point at the image centre).

\subsection{Baseline}
\label{sec:baseline}

On \code{swiss\_knife.mp4} (628 frames, $1280\times720$ HEVC, handheld orbit):

\begin{center}
\begin{tabular}{ll}
\toprule
\textbf{Metric} & \textbf{Value}\\
\midrule
Registered views & 22 / 63\\
Sparse points / RMSE & 1{,}175 / 2.086\,px\\
Dense depth vs sparse & median 4.8\pct, p90 44.8\pct (13{,}659 samples)\\
Raw mesh & 109{,}933 faces, 525 patches\\
Final mesh & 20{,}000 faces, 36 patches\\
Mesh topology & open (partial scan), Euler $=-102$\\
Wall time & $\approx$55\,s (decode 23.5\,s, dense $\approx$25\,s)\\
\bottomrule
\end{tabular}
\end{center}

\subsection{Accuracy pass}
\label{sec:accuracy}

\begin{center}
\begin{tabular}{lccc}
\toprule
\textbf{Capture} & \textbf{Depth median / p90} & \textbf{Patches} & \textbf{Euler}\\
\midrule
\code{swiss\_knife} (specular) & 4.8$\to$\textbf{4.0}\pct / 44.8$\to$\textbf{29.4}\pct & 34$\to$33 & $-104\to-97$\\
\code{lip\_balm} 18--27\,s (white) & 19.0$\to$\textbf{14.8}\pct / 77$\to$79\pct & 41 & $-196$\\
\code{keyboard} 15--25\,s (textured) & 6.8$\to$\textbf{4.5}\pct / 55.3$\to$52.8\pct & 2$\to$4 & $-36\to-9$\\
\bottomrule
\end{tabular}
\end{center}

The improvements came from parabolic sub-pixel depth refinement, 8-path SGM, a depth speckle
filter, and feature-preserving bilateral mesh smoothing.

\subsection{Presets}
\label{sec:presets}

The four Studio presets are complete configs (they replace, not merge); \code{fast},
\code{balanced} and \code{quality} pin explicit dense parameters, while \code{auto} derives them
from the measured frame texture.

\begin{center}
\begin{tabular}{lccccc}
\toprule
\textbf{Preset} & \textbf{step} & \textbf{scale} & \textbf{planes} & \textbf{voxel} & \textbf{faces}\\
\midrule
fast & 20 & 0.35 & 64 & 192 & 20{,}000\\
balanced & 10 & 0.50 & 96 & 256 & 100{,}000\\
quality & 5 & 0.75 & 128 & 320 & 80{,}000\\
auto & 10 & 0.50 & 96 & 256 & 100{,}000\\
\bottomrule
\end{tabular}
\end{center}

Measured on \code{swiss\_knife.mp4}:

\begin{center}
\begin{tabular}{lcccc}
\toprule
\textbf{Preset} & \textbf{Registered} & \textbf{RMSE} & \textbf{Depth} & \textbf{Time}\\
\midrule
fast & 8/32 & 1.58\,px & 2.0\pct & 32\,s\\
balanced & 19/63 & 2.09\,px & 3.9\pct & 77\,s\\
quality & 40/126 & 2.44\,px & 3.3\pct & 340\,s\\
auto & 19/63 & 2.09\,px & 3.9\pct & 76\,s\\
\bottomrule
\end{tabular}
\end{center}

\subsection{Reference captures}
\label{sec:captures}

\begin{center}
\begin{tabular}{lcccl}
\toprule
\textbf{Clip} & \textbf{Resolution} & \textbf{Views} & \textbf{Depth} & \textbf{Result}\\
\midrule
\code{swiss\_knife} & $1280\times720$ & $\sim$22/63 & 4.0\pct & 33 patches\\
\code{keyboard} & $480\times854$ & 63 & 4.5\pct & 4 patches\\
\code{stand} & $480\times854$ & 140/380 & 2.0\pct & 20 patches\\
\code{lip\_balm} & $480\times854$ & 66 & 14.8\pct & 41 patches\\
\code{remote} & $854\times480$ & 20/151 & 4.2\pct med / 56\pct p90 & near-empty\\
\code{tube} & $854\times480$ & 2/91 & --- & fails (27 kp/frame)\\
\bottomrule
\end{tabular}
\end{center}

\subsection{Foreground isolation}
\label{sec:iso-res}

With \code{--isolate-object} on \code{swiss\_knife.mp4} (identical views and RMSE), the mesh went
from \textbf{49 patches / Euler $-345$} to \textbf{1 patch / Euler $-64$}, by cropping the dense
cloud to the object box and rescaling the TSDF resolution.

\subsection{Determinism and caching}
\label{sec:det}

Fresh runs, cache hits and cache misses are byte-identical. The sparse-stage cache (decoded
frames + tracks, keyed by video SHA-256 and stage parameters) reduced an 11\,s clip's sparse
stage from 42.7\,s to 13.2\,s.

\section{Trust and measurement}
\label{sec:trust}

No single number proves a mesh is good, so the report triangulates several: sparse reprojection
RMSE (camera/point solve), depth-vs-sparse error (a low median with a high p90 means good
typical depth and a bad outlier tail), mesh topology (patch count, Euler characteristic,
boundary ratio), and a self-assessment confidence. The self-assessment is a consistency signal,
not a guarantee: it has been observed to report ``high'' on a run whose mesh covered almost no
surface, which is why coverage and continuity are treated separately.

Thirteen named failure codes (e.g.\ \code{low-registration}, \code{empty-mesh},
\code{calibration-failed}) carry actionable hints, so the UI shows a specific message rather than
a stack trace.

\section{Iteration history}
\label{sec:history}

Each entry is a real failure with symptom, root cause, fix and measured before/after.

\begin{enumerate}[leftmargin=1.5em]
\item \textbf{Blobs and holes.} Winner-take-all NCC had no spatial smoothness and depth was
quantised. Fixed with sub-pixel refinement, 8-path SGM, a speckle filter and bilateral
smoothing (depth median $4.8\to4.0\pct$).
\item \textbf{``Not enough views'' on landscape clips.} New clips were $854\times480$ landscape
but the intrinsics were from $480\times854$ portrait, so $c_x,c_y$ were swapped. Registration
went from 12/151 to 24/150, and 84/84 on the trimmed segment.
\item \textbf{Low keypoint density.} Blurry or textureless captures yield $\sim$30 keypoints per
frame; the tool warns below 80. The tube clip measured 27 keypoints/frame $\to$ 2/91 views.
\item \textbf{Non-deterministic cache.} \code{serde\_json} round-trips f32-promoted f64
coordinates one ULP off, and the frame cache stored gray/RGB in blocks but read them
interleaved. Fixed by storing raw f64 bits and matching the block layout; reruns are now
byte-identical.
\item \textbf{Preset drift.} Presets were partial overrides (and \code{balanced} was empty), so
they inherited the previous mode. Fixed by making presets complete configs that replace the
current one.
\item \textbf{Fragile metric scale.} kornia's \code{tag\_scale} triangulates from one
widest-baseline pair; on an orbit the tag anchored a 10\,cm tag to a $\sim$20--66\,m mesh. The
application now uses N-view DLT with a plausibility warning, and the root cause (gauge drift
without loop closure) is filed upstream as \cite{loopclosure}.
\item \textbf{Self-assessment over-claim.} A run reported 0.77 confidence for a 768-face mesh;
the gap is documented and coverage is treated separately.
\item \textbf{Whole-scene fusion.} Fusing the whole room shatters the object; foreground
isolation + adaptive resolution fixed it (49$\to$1 patch).
\end{enumerate}

\section{Limitations and future work}
\label{sec:future}

\begin{itemize}[leftmargin=1.5em]
\item \textbf{Scale drift} on long orbits without loop closure can make a ``metric'' map wrong
even when a tag anchors it; a plausibility warning is the current guard, and the fix (pose-graph
loop closure) is proposed upstream \cite{loopclosure}.
\item \textbf{Foreground isolation} is heuristic and centred-object specific; it declines on
clips with too few central points rather than guessing.
\item \textbf{Non-watertight meshes.} Marching tetrahedra does not guarantee a closed surface;
screened Poisson \cite{kazhdan2013} is the planned route to watertightness.
\item \textbf{Textureless / specular / dark} surfaces remain the dominant failure mode.
\end{itemize}

Further upstream improvements (track merging, feature-path covariance, optimal triangulation,
adaptive thresholds) are tracked in the repository's \code{docs/upstream-kornia-issues.md}.

\begin{thebibliography}{99}

\bibitem{lowe2004} D.~G. Lowe. Distinctive image features from scale-invariant keypoints.
\emph{International Journal of Computer Vision}, 2004.

\bibitem{rublee2011} E.~Rublee, V.~Rabaud, K.~Konolige, G.~Bradski. ORB: an efficient
alternative to SIFT or SURF. \emph{ICCV}, 2011.

\bibitem{hartley2004} R.~Hartley, A.~Zisserman. \emph{Multiple View Geometry in Computer
Vision}. Cambridge University Press, 2004.

\bibitem{fischler1981} M.~A. Fischler, R.~C. Bolles. Random sample consensus: a paradigm for
model fitting. \emph{Communications of the ACM}, 1981.

\bibitem{lepetit2009} V.~Lepetit, F.~Moreno-Noguer, P.~Fua. EPnP: an accurate $O(n)$ solution
to the PnP problem. \emph{International Journal of Computer Vision}, 2009.

\bibitem{wald1945} A.~Wald. Sequential tests of statistical hypotheses. \emph{Annals of
Mathematical Statistics}, 1945.

\bibitem{triggs2000} B.~Triggs, P.~McLauchlan, R.~Hartley, A.~Fitzgibbon. Bundle adjustment ---
a modern synthesis. \emph{Vision Algorithms: Theory and Practice}, 2000.

\bibitem{lourakis2009} M.~Lourakis, A.~Argyros. SBA: a software package for generic sparse
bundle adjustment. \emph{ACM Transactions on Mathematical Software}, 2009.

\bibitem{zhang2000} Z.~Zhang. A flexible new technique for camera calibration. \emph{IEEE
Transactions on Pattern Analysis and Machine Intelligence}, 2000.

\bibitem{olson2011} E.~Olson. AprilTag: a robust and flexible visual fiducial system.
\emph{ICRA}, 2011.

\bibitem{schonberger2016} J.~L. Sch\"onberger, J.-M. Frahm. Structure-from-motion revisited.
\emph{CVPR}, 2016.

\bibitem{collins1996} R.~T. Collins. A space-sweep approach to true multi-image matching.
\emph{CVPR}, 1996.

\bibitem{zabih1994} R.~Zabih, J.~Woodfill. Non-parametric local transforms for computing visual
correspondence. \emph{ECCV}, 1994.

\bibitem{hirschmuller2008} H.~Hirschm\"uller. Stereo processing by semi-global matching and
mutual information. \emph{IEEE Transactions on Pattern Analysis and Machine Intelligence}, 2008.

\bibitem{curless1996} B.~Curless, M.~Levoy. A volumetric method for building complex models
from range images. \emph{SIGGRAPH}, 1996.

\bibitem{newcombe2011} R.~A. Newcombe et al. KinectFusion: real-time dense surface mapping and
tracking. \emph{ISMAR}, 2011.

\bibitem{lorensen1987} W.~E. Lorensen, H.~E. Cline. Marching cubes: a high resolution 3D surface
construction algorithm. \emph{SIGGRAPH}, 1987.

\bibitem{taubin1995} G.~Taubin. Curve and surface smoothing without shrinkage. \emph{ICCV}, 1995.

\bibitem{garland1997} M.~Garland, P.~S. Heckbert. Surface simplification using quadric error
metrics. \emph{SIGGRAPH}, 1997.

\bibitem{kazhdan2013} M.~Kazhdan, H.~Hoppe. Screened Poisson surface reconstruction. \emph{ACM
Transactions on Graphics}, 2013.

\bibitem{huber1964} P.~J. Huber. Robust estimation of a location parameter. \emph{Annals of
Mathematical Statistics}, 1964.

\bibitem{umeyama1991} S.~Umeyama. Least-squares estimation of transformation parameters between
two point patterns. \emph{IEEE Transactions on Pattern Analysis and Machine Intelligence}, 1991.

\bibitem{loopclosure} \code{meshprod} issue: Loop closure / pose-graph optimization for
incremental SfM. \url{https://github.com/kornia/kornia-rs/issues/1365}.

\end{thebibliography}

\end{document}
